\textbf{Proof.}
Apply \(\mathrm{thm:global-elementary-minimax-bracket}\).  Its upper bound is already the asserted one.  For the lower bound, choose the range constant \(\rho_\epsilon>0\) no larger than \(1/4\).  If \(d\le\rho_\epsilon n\log(en)\), then
\[
 \frac d{n\log(en)}\le\rho_\epsilon,
 \qquad
 \frac{\sqrt d}{n}
 \le\sqrt{\rho_\epsilon\frac{\log(en)}n}
 \le\sqrt{\rho_\epsilon},
\]
where \(\log(en)/n\le1\) for every integer \(n\ge1\).  Hence the sum inside the global truncation is at most \(\sqrt{\rho_\epsilon}+\rho_\epsilon<1\), so the truncation disappears.  This gives the displayed in-range lower bound after retaining the same \(c_\epsilon\).

The comparison
\[
 \frac{\sqrt d/n}{d/\{n\log(en)\}}
 =\frac{\log(en)}{\sqrt d}
\]
proves the elbow and dominance assertions.  The necessary and sufficient consistency statements, the parametric iff, and the explicit description of the remaining gap are restrictions of the corresponding conclusions of the global theorem to the displayed range.  No additional support, positivity, or boundary condition is introduced by this restriction.